% Options for packages loaded elsewhere
\PassOptionsToPackage{unicode}{hyperref}
\PassOptionsToPackage{hyphens}{url}
%
\documentclass[
]{article}
\usepackage{amsmath,amssymb}
\usepackage{iftex}
\ifPDFTeX
  \usepackage[T1]{fontenc}
  \usepackage[utf8]{inputenc}
  \usepackage{textcomp} % provide euro and other symbols
\else % if luatex or xetex
  \usepackage{unicode-math} % this also loads fontspec
  \defaultfontfeatures{Scale=MatchLowercase}
  \defaultfontfeatures[\rmfamily]{Ligatures=TeX,Scale=1}
\fi
\usepackage{lmodern}
\ifPDFTeX\else
  % xetex/luatex font selection
\fi
% Use upquote if available, for straight quotes in verbatim environments
\IfFileExists{upquote.sty}{\usepackage{upquote}}{}
\IfFileExists{microtype.sty}{% use microtype if available
  \usepackage[]{microtype}
  \UseMicrotypeSet[protrusion]{basicmath} % disable protrusion for tt fonts
}{}
\makeatletter
\@ifundefined{KOMAClassName}{% if non-KOMA class
  \IfFileExists{parskip.sty}{%
    \usepackage{parskip}
  }{% else
    \setlength{\parindent}{0pt}
    \setlength{\parskip}{6pt plus 2pt minus 1pt}}
}{% if KOMA class
  \KOMAoptions{parskip=half}}
\makeatother
\usepackage{xcolor}
\usepackage[margin=1in]{geometry}
\usepackage{color}
\usepackage{fancyvrb}
\newcommand{\VerbBar}{|}
\newcommand{\VERB}{\Verb[commandchars=\\\{\}]}
\DefineVerbatimEnvironment{Highlighting}{Verbatim}{commandchars=\\\{\}}
% Add ',fontsize=\small' for more characters per line
\usepackage{framed}
\definecolor{shadecolor}{RGB}{248,248,248}
\newenvironment{Shaded}{\begin{snugshade}}{\end{snugshade}}
\newcommand{\AlertTok}[1]{\textcolor[rgb]{0.94,0.16,0.16}{#1}}
\newcommand{\AnnotationTok}[1]{\textcolor[rgb]{0.56,0.35,0.01}{\textbf{\textit{#1}}}}
\newcommand{\AttributeTok}[1]{\textcolor[rgb]{0.13,0.29,0.53}{#1}}
\newcommand{\BaseNTok}[1]{\textcolor[rgb]{0.00,0.00,0.81}{#1}}
\newcommand{\BuiltInTok}[1]{#1}
\newcommand{\CharTok}[1]{\textcolor[rgb]{0.31,0.60,0.02}{#1}}
\newcommand{\CommentTok}[1]{\textcolor[rgb]{0.56,0.35,0.01}{\textit{#1}}}
\newcommand{\CommentVarTok}[1]{\textcolor[rgb]{0.56,0.35,0.01}{\textbf{\textit{#1}}}}
\newcommand{\ConstantTok}[1]{\textcolor[rgb]{0.56,0.35,0.01}{#1}}
\newcommand{\ControlFlowTok}[1]{\textcolor[rgb]{0.13,0.29,0.53}{\textbf{#1}}}
\newcommand{\DataTypeTok}[1]{\textcolor[rgb]{0.13,0.29,0.53}{#1}}
\newcommand{\DecValTok}[1]{\textcolor[rgb]{0.00,0.00,0.81}{#1}}
\newcommand{\DocumentationTok}[1]{\textcolor[rgb]{0.56,0.35,0.01}{\textbf{\textit{#1}}}}
\newcommand{\ErrorTok}[1]{\textcolor[rgb]{0.64,0.00,0.00}{\textbf{#1}}}
\newcommand{\ExtensionTok}[1]{#1}
\newcommand{\FloatTok}[1]{\textcolor[rgb]{0.00,0.00,0.81}{#1}}
\newcommand{\FunctionTok}[1]{\textcolor[rgb]{0.13,0.29,0.53}{\textbf{#1}}}
\newcommand{\ImportTok}[1]{#1}
\newcommand{\InformationTok}[1]{\textcolor[rgb]{0.56,0.35,0.01}{\textbf{\textit{#1}}}}
\newcommand{\KeywordTok}[1]{\textcolor[rgb]{0.13,0.29,0.53}{\textbf{#1}}}
\newcommand{\NormalTok}[1]{#1}
\newcommand{\OperatorTok}[1]{\textcolor[rgb]{0.81,0.36,0.00}{\textbf{#1}}}
\newcommand{\OtherTok}[1]{\textcolor[rgb]{0.56,0.35,0.01}{#1}}
\newcommand{\PreprocessorTok}[1]{\textcolor[rgb]{0.56,0.35,0.01}{\textit{#1}}}
\newcommand{\RegionMarkerTok}[1]{#1}
\newcommand{\SpecialCharTok}[1]{\textcolor[rgb]{0.81,0.36,0.00}{\textbf{#1}}}
\newcommand{\SpecialStringTok}[1]{\textcolor[rgb]{0.31,0.60,0.02}{#1}}
\newcommand{\StringTok}[1]{\textcolor[rgb]{0.31,0.60,0.02}{#1}}
\newcommand{\VariableTok}[1]{\textcolor[rgb]{0.00,0.00,0.00}{#1}}
\newcommand{\VerbatimStringTok}[1]{\textcolor[rgb]{0.31,0.60,0.02}{#1}}
\newcommand{\WarningTok}[1]{\textcolor[rgb]{0.56,0.35,0.01}{\textbf{\textit{#1}}}}
\usepackage{graphicx}
\makeatletter
\def\maxwidth{\ifdim\Gin@nat@width>\linewidth\linewidth\else\Gin@nat@width\fi}
\def\maxheight{\ifdim\Gin@nat@height>\textheight\textheight\else\Gin@nat@height\fi}
\makeatother
% Scale images if necessary, so that they will not overflow the page
% margins by default, and it is still possible to overwrite the defaults
% using explicit options in \includegraphics[width, height, ...]{}
\setkeys{Gin}{width=\maxwidth,height=\maxheight,keepaspectratio}
% Set default figure placement to htbp
\makeatletter
\def\fps@figure{htbp}
\makeatother
\setlength{\emergencystretch}{3em} % prevent overfull lines
\providecommand{\tightlist}{%
  \setlength{\itemsep}{0pt}\setlength{\parskip}{0pt}}
\setcounter{secnumdepth}{-\maxdimen} % remove section numbering
\ifLuaTeX
  \usepackage{selnolig}  % disable illegal ligatures
\fi
\usepackage{bookmark}
\IfFileExists{xurl.sty}{\usepackage{xurl}}{} % add URL line breaks if available
\urlstyle{same}
\hypersetup{
  pdftitle={Statistical Method II Assignment II},
  pdfauthor={Daniel Oluwafemi Olofin},
  hidelinks,
  pdfcreator={LaTeX via pandoc}}

\title{Statistical Method II Assignment II}
\author{Daniel Oluwafemi Olofin}
\date{2025-04-05}

\begin{document}
\maketitle

\begin{Shaded}
\begin{Highlighting}[]
\FunctionTok{library}\NormalTok{(readr)}
\end{Highlighting}
\end{Shaded}

\begin{verbatim}
## Warning: package 'readr' was built under R version 4.2.3
\end{verbatim}

\begin{Shaded}
\begin{Highlighting}[]
\FunctionTok{library}\NormalTok{(tidyverse)}
\end{Highlighting}
\end{Shaded}

\begin{verbatim}
## Warning: package 'tidyverse' was built under R version 4.2.3
\end{verbatim}

\begin{verbatim}
## Warning: package 'ggplot2' was built under R version 4.2.3
\end{verbatim}

\begin{verbatim}
## Warning: package 'tibble' was built under R version 4.2.3
\end{verbatim}

\begin{verbatim}
## Warning: package 'tidyr' was built under R version 4.2.3
\end{verbatim}

\begin{verbatim}
## Warning: package 'purrr' was built under R version 4.2.3
\end{verbatim}

\begin{verbatim}
## Warning: package 'dplyr' was built under R version 4.2.3
\end{verbatim}

\begin{verbatim}
## Warning: package 'stringr' was built under R version 4.2.3
\end{verbatim}

\begin{verbatim}
## Warning: package 'forcats' was built under R version 4.2.3
\end{verbatim}

\begin{verbatim}
## Warning: package 'lubridate' was built under R version 4.2.3
\end{verbatim}

\begin{verbatim}
## -- Attaching core tidyverse packages ------------------------ tidyverse 2.0.0 --
## v dplyr     1.1.4     v purrr     1.0.2
## v forcats   1.0.0     v stringr   1.5.1
## v ggplot2   3.5.1     v tibble    3.2.1
## v lubridate 1.9.3     v tidyr     1.3.1
## -- Conflicts ------------------------------------------ tidyverse_conflicts() --
## x dplyr::filter() masks stats::filter()
## x dplyr::lag()    masks stats::lag()
## i Use the conflicted package (<http://conflicted.r-lib.org/>) to force all conflicts to become errors
\end{verbatim}

\begin{Shaded}
\begin{Highlighting}[]
\NormalTok{House\_data }\OtherTok{\textless{}{-}} \FunctionTok{read\_table}\NormalTok{(}\StringTok{"House\_data.txt"}\NormalTok{)}
\end{Highlighting}
\end{Shaded}

\begin{verbatim}
## Warning: Missing column names filled in: 'X8' [8]
\end{verbatim}

\begin{verbatim}
## 
## -- Column specification --------------------------------------------------------
## cols(
##   POP = col_double(),
##   POPCH = col_double(),
##   CHILD = col_double(),
##   LUNCH = col_double(),
##   INCOME = col_double(),
##   CRIME = col_double(),
##   CRIMECH = col_double(),
##   X8 = col_logical()
## )
\end{verbatim}

\begin{verbatim}
## Warning: 42 parsing failures.
## row col  expected    actual             file
##   1  -- 8 columns 7 columns 'House_data.txt'
##   2  -- 8 columns 7 columns 'House_data.txt'
##   3  -- 8 columns 7 columns 'House_data.txt'
##   4  -- 8 columns 7 columns 'House_data.txt'
##   5  -- 8 columns 7 columns 'House_data.txt'
## ... ... ......... ......... ................
## See problems(...) for more details.
\end{verbatim}

\begin{Shaded}
\begin{Highlighting}[]
\FunctionTok{View}\NormalTok{(House\_data)}

\CommentTok{\#Data Managemnet}
\NormalTok{House\_data }\SpecialCharTok{\%\textgreater{}\%} 
  \FunctionTok{select}\NormalTok{(}\SpecialCharTok{{-}}\DecValTok{8}\NormalTok{) }\OtherTok{{-}\textgreater{}}\NormalTok{ House\_data}

\FunctionTok{str\_to\_sentence}\NormalTok{(}\FunctionTok{names}\NormalTok{(House\_data)) }\OtherTok{{-}\textgreater{}} \FunctionTok{names}\NormalTok{(House\_data)}
\end{Highlighting}
\end{Shaded}

\section{Question 2}\label{question-2}

\subsection{1. Show the linear regression output in
R.}\label{show-the-linear-regression-output-in-r.}

\begin{Shaded}
\begin{Highlighting}[]
\NormalTok{Crime\_model }\OtherTok{\textless{}{-}} \FunctionTok{lm}\NormalTok{(Crime }\SpecialCharTok{\textasciitilde{}}\NormalTok{., }\AttributeTok{data =}\NormalTok{ House\_data)}

\FunctionTok{summary}\NormalTok{(Crime\_model)}
\end{Highlighting}
\end{Shaded}

\begin{verbatim}
## 
## Call:
## lm(formula = Crime ~ ., data = House_data)
## 
## Residuals:
##      Min       1Q   Median       3Q      Max 
## -117.010  -37.285   -2.292   28.046  241.633 
## 
## Coefficients:
##             Estimate Std. Error t value Pr(>|t|)    
## (Intercept)  67.5618    58.8035   1.149   0.2580    
## Pop          -5.2679     3.1061  -1.696   0.0983 .  
## Popch         1.2279     0.9558   1.285   0.2069    
## Child        -7.0990     1.6171  -4.390 9.11e-05 ***
## Lunch         2.8437     0.5304   5.361 4.60e-06 ***
## Income        4.5877     1.6586   2.766   0.0088 ** 
## Crimech       0.5862     0.7550   0.776   0.4424    
## ---
## Signif. codes:  0 '***' 0.001 '**' 0.01 '*' 0.05 '.' 0.1 ' ' 1
## 
## Residual standard error: 66.18 on 37 degrees of freedom
## Multiple R-squared:  0.658,  Adjusted R-squared:  0.6025 
## F-statistic: 11.86 on 6 and 37 DF,  p-value: 2.187e-07
\end{verbatim}

\subsection{Construct the ANOVA table and provide interpretation for
each single number in the
table.}\label{construct-the-anova-table-and-provide-interpretation-for-each-single-number-in-the-table.}

\begin{Shaded}
\begin{Highlighting}[]
\FunctionTok{anova}\NormalTok{(Crime\_model)}
\end{Highlighting}
\end{Shaded}

\begin{verbatim}
## Analysis of Variance Table
## 
## Response: Crime
##           Df Sum Sq Mean Sq F value    Pr(>F)    
## Pop        1  61486   61486 14.0394 0.0006095 ***
## Popch      1  19407   19407  4.4313 0.0421361 *  
## Child      1  30561   30561  6.9780 0.0120193 *  
## Lunch      1 165490  165490 37.7870  3.97e-07 ***
## Income     1  32178   32178  7.3474 0.0101210 *  
## Crimech    1   2641    2641  0.6029 0.4423983    
## Residuals 37 162043    4380                      
## ---
## Signif. codes:  0 '***' 0.001 '**' 0.01 '*' 0.05 '.' 0.1 ' ' 1
\end{verbatim}

\subsection{3. Perform an overall F-test of whether there is a
regression relation between CRIME Y and the set of variables in the
model, using α =
0.05.}\label{perform-an-overall-f-test-of-whether-there-is-a-regression-relation-between-crime-y-and-the-set-of-variables-in-the-model-using-ux3b1-0.05.}

\subsubsection{\texorpdfstring{3.1. .State the null and alternative
hypotheses \textbf{for Overall
F-test}.}{3.1. .State the null and alternative hypotheses for Overall F-test.}}\label{state-the-null-and-alternative-hypotheses-for-overall-f-test.}

\textbf{Null Hypothesis (\(H_0\)):} \(\beta_1 = 0\) for i = 1,2\ldots, 6
(None of the predictors have a linear relationship with the Crime
variable.)

\textbf{Alternative Hypothesis (\(H_a\)):} \(\beta_1 \neq 0\) for i =
1,2\ldots, 6 (At least one predictor has a statistically significant
effect on the Crime variable.)

\subsubsection{3.2. What is the value of the test statistic for this
test?}\label{what-is-the-value-of-the-test-statistic-for-this-test}

The test statistic is \textbf{F-statistic (11.86)}.

\subsubsection{3.3. What degrees of freedom are used for this
test?}\label{what-degrees-of-freedom-are-used-for-this-test}

The degrees of freedom for the overall model is \textbf{37}.

\subsubsection{3.4. What is the p-value for this
test?}\label{what-is-the-p-value-for-this-test}

The associated p-value for the overall F-statistic is
\textbf{2.187e-07}.

\subsubsection{3.5. Should the null hypothesis be rejected?
Explain.}\label{should-the-null-hypothesis-be-rejected-explain.}

Since the P-value is \textless{} 0.05 level of significance, we have
enough evidence to reject the null hypothesis.

\subsubsection{3.6. State your
conclusions.}\label{state-your-conclusions.}

we conclude that at least one predictor has a statistically significant
effect on the Crime variable.

\subsubsection{3.1. State the null and alternative hypotheses for
Pop.}\label{state-the-null-and-alternative-hypotheses-for-pop.}

\textbf{Null Hypothesis (\(H_0\)):} \(\beta_1 = 0\) (No linear
relationship)

\textbf{Alternative Hypothesis (\(H_a\)):} \(\beta_1 \neq 0\) (Linear
relationship exists)

\subsubsection{3.2. What is the value of the test statistic for this
test?}\label{what-is-the-value-of-the-test-statistic-for-this-test-1}

The test statistic is \textbf{t-statistic (-1.696)}.

\subsubsection{3.3. What degrees of freedom are used for this
test?}\label{what-degrees-of-freedom-are-used-for-this-test-1}

The degrees of freedom for \(\beta_1 = 37\).

\subsubsection{3.4. What is the p-value for this
test?}\label{what-is-the-p-value-for-this-test-1}

The associated p-value for the overall t-statistic is \textbf{0.0983}.

\subsubsection{3.5. Should the null hypothesis be rejected?
Explain.}\label{should-the-null-hypothesis-be-rejected-explain.-1}

Since the P-value is \textgreater{} 0.05 level of significance, we have
enough evidence to fail to reject the null hypothesis.

\subsubsection{3.6. State your
conclusions.}\label{state-your-conclusions.-1}

We conclude that Pop doesn't have a statistically significant effect on
the Crime variable.

\subsubsection{3.1. State the null and alternative hypotheses for
Popch.}\label{state-the-null-and-alternative-hypotheses-for-popch.}

\textbf{Null Hypothesis (\(H_0\)):} \(\beta_2 = 0\) (No linear
relationship)

\textbf{Alternative Hypothesis (\(H_a\)):} \(\beta_2 \neq 0\) (Linear
relationship exists)

\subsubsection{3.2. What is the value of the test statistic for this
test?}\label{what-is-the-value-of-the-test-statistic-for-this-test-2}

The test statistic is \textbf{t-statistic (1.285)}.

\subsubsection{3.3. What degrees of freedom are used for this
test?}\label{what-degrees-of-freedom-are-used-for-this-test-2}

The degrees of freedom for \(\beta_2 = 37\).

\subsubsection{3.4. What is the p-value for this
test?}\label{what-is-the-p-value-for-this-test-2}

The associated p-value for the overall t-statistic is \textbf{0.2580}.

\subsubsection{3.5. Should the null hypothesis be rejected?
Explain.}\label{should-the-null-hypothesis-be-rejected-explain.-2}

Since the P-value is \textgreater{} 0.05 level of significance, we have
enough evidence to fail to reject the null hypothesis.

\subsubsection{3.6. State your
conclusions.}\label{state-your-conclusions.-2}

We conclude that Popch doesn't have a statistically significant effect
on the Crime variable.

\subsubsection{3.1. State the null and alternative hypotheses for
child.}\label{state-the-null-and-alternative-hypotheses-for-child.}

\textbf{Null Hypothesis (\(H_0\)):} \(\beta_3 = 0\) (No linear
relationship)

\textbf{Alternative Hypothesis (\(H_a\)):} \(\beta_3 \neq 0\) (Linear
relationship exists)

\subsubsection{3.2. What is the value of the test statistic for this
test?}\label{what-is-the-value-of-the-test-statistic-for-this-test-3}

The test statistic is \textbf{t-statistic (-4.390)}.

\subsubsection{3.3. What degrees of freedom are used for this
test?}\label{what-degrees-of-freedom-are-used-for-this-test-3}

The degrees of freedom for \(\beta_3 = 37\).

\subsubsection{3.4. What is the p-value for this
test?}\label{what-is-the-p-value-for-this-test-3}

The associated p-value for the overall t-statistic is \textbf{9.11e-05}.

\subsubsection{3.5. Should the null hypothesis be rejected?
Explain.}\label{should-the-null-hypothesis-be-rejected-explain.-3}

Since the P-value is \textless{} 0.05 level of significance, we have
enough evidence to reject the null hypothesis.

\subsubsection{3.6. State your
conclusions.}\label{state-your-conclusions.-3}

We conclude that Child has a statistically significant effect on the
Crime variable.

\subsubsection{3.1. State the null and alternative hypotheses for
Lunch.}\label{state-the-null-and-alternative-hypotheses-for-lunch.}

\textbf{Null Hypothesis (\(H_0\)):} \(\beta_4 = 0\) (No linear
relationship)

\textbf{Alternative Hypothesis (\(H_a\)):} \(\beta_4 \neq 0\) (Linear
relationship exists)

\subsubsection{3.2. What is the value of the test statistic for this
test?}\label{what-is-the-value-of-the-test-statistic-for-this-test-4}

The test statistic is \textbf{t-statistic (5.361)}.

\subsubsection{3.3. What degrees of freedom are used for this
test?}\label{what-degrees-of-freedom-are-used-for-this-test-4}

The degrees of freedom for \(\beta_4 = 37\).

\subsubsection{3.4. What is the p-value for this
test?}\label{what-is-the-p-value-for-this-test-4}

The associated p-value for the overall t-statistic is \textbf{4.60e-06}.

\subsubsection{3.5. Should the null hypothesis be rejected?
Explain.}\label{should-the-null-hypothesis-be-rejected-explain.-4}

Since the P-value is \textless{} 0.05 level of significance, we have
enough evidence to reject the null hypothesis.

\subsubsection{3.6. State your
conclusions.}\label{state-your-conclusions.-4}

We conclude that Lunch has a statistically significant effect on the
Crime variable.

\subsubsection{3.1. State the null and alternative hypotheses for
Income.}\label{state-the-null-and-alternative-hypotheses-for-income.}

\textbf{Null Hypothesis (\(H_0\)):} \(\beta_5 = 0\) (No linear
relationship)

\textbf{Alternative Hypothesis (\(H_a\)):} \(\beta_5 \neq 0\) (Linear
relationship exists)

\subsubsection{3.2. What is the value of the test statistic for this
test?}\label{what-is-the-value-of-the-test-statistic-for-this-test-5}

The test statistic is \textbf{t-statistic (2.766)}.

\subsubsection{3.3. What degrees of freedom are used for this
test?}\label{what-degrees-of-freedom-are-used-for-this-test-5}

The degrees of freedom for \(\beta_5 = 37\).

\subsubsection{3.4. What is the p-value for this
test?}\label{what-is-the-p-value-for-this-test-5}

The associated p-value for the overall t-statistic is \textbf{0.0088}.

\subsubsection{3.5. Should the null hypothesis be rejected?
Explain.}\label{should-the-null-hypothesis-be-rejected-explain.-5}

Since the P-value is \textless{} 0.05 level of significance, we have
enough evidence to reject the null hypothesis.

\subsubsection{3.6. State your
conclusions.}\label{state-your-conclusions.-5}

We conclude that Income has a statistically significant effect on the
Crime variable.

\subsubsection{3.1. State the null and alternative hypotheses for
child.}\label{state-the-null-and-alternative-hypotheses-for-child.-1}

\textbf{Null Hypothesis (\(H_0\)):} \(\beta_6 = 0\) (No linear
relationship)

\textbf{Alternative Hypothesis (\(H_a\)):} \(\beta_6 \neq 0\) (Linear
relationship exists)

\subsubsection{3.2. What is the value of the test statistic for this
test?}\label{what-is-the-value-of-the-test-statistic-for-this-test-6}

The test statistic is \textbf{t-statistic (0.776)}.

\subsubsection{3.3. What degrees of freedom are used for this
test?}\label{what-degrees-of-freedom-are-used-for-this-test-6}

The degrees of freedom for \(\beta_6 = 37\).

\subsubsection{3.4. What is the p-value for this
test?}\label{what-is-the-p-value-for-this-test-6}

The associated p-value for the overall t-statistic is \textbf{0.4424}.

\subsubsection{3.5. Should the null hypothesis be rejected?
Explain.}\label{should-the-null-hypothesis-be-rejected-explain.-6}

Since the P-value is \textless{} 0.05 level of significance, we have
enough evidence to reject the null hypothesis.

\subsubsection{3.6. State your
conclusions.}\label{state-your-conclusions.-6}

We conclude that Crimech doesn't have a statistically significant effect
on the Crime variable.

\subsection{4. Test simultaneously whether 𝑃𝑂𝑃𝐶𝐻 and 𝐶𝑅𝐼𝑀𝐸𝐶𝐻 must be
removed from the model. Test full vs reduced model. (23pts, 3 each, only
the first one is
2)}\label{test-simultaneously-whether-ux1d443ux1d442ux1d443ux1d436ux1d43b-and-ux1d436ux1d445ux1d43cux1d440ux1d438ux1d436ux1d43b-must-be-removed-from-the-model.-test-full-vs-reduced-model.-23pts-3-each-only-the-first-one-is-2}

\end{document}
